\section*{Capítulo \ref{ch:b3:remote-sensing-gis} --- Teledetección y SIG}

\begin{solution}{exo:b3:remote-sensing-gis:1}
Campo: $(0.42 - 0.06)/(0.42 + 0.06) = 0.75$, vegetación verde densa, un cultivo
en plena foliación. Campo arado: $(0.28 - 0.21)/(0.28 + 0.21) = 0.14$, casi
suelo desnudo, un suelo que se aclara ligeramente hacia el infrarrojo.
\end{solution}

\begin{solution}{exo:b3:remote-sensing-gis:2}
(a)~$185000/30 \approx 6170$ píxeles; $185000/15 \approx 12300$ en la
banda pancromática. (b)~$2^{12} = 4096$ niveles; $2^8 = 256$. (c)~16 días
por sí solo; con Landsat~9 en la misma órbita y desfasado 8 días, se reduce a 8 días.
\end{solution}

\begin{solution}{exo:b3:remote-sensing-gis:3}
(a)~Una franja equivale a $\lambda/2 = \qty{27.7}{mm}$; seis franjas suponen
\qty{166}{mm} de variación en la \omterm{def:b3:remote-sensing-gis:insar}{línea de visión}, aquí alejándose del satélite, ya que
la cuenca experimenta subsidencia. (b)~$56.6/2 = \qty{28.3}{mm}$.
\end{solution}

\begin{solution}{exo:b3:remote-sensing-gis:4}
(a)~$t = 2d/c = 2\times\num{5e5}/\num{3.00e8} = \qty{3.34}{ms}$.
(b)~$c\,\delta t/2 = \num{3.00e8}\times\num{8e-10}/2 = \qty{0.12}{m}$; una
superficie muestreada por muchos fotones se localiza mucho mejor, la precisión
disminuye como la inversa de la raíz cuadrada de su número.
\end{solution}

\begin{solution}{exo:b3:remote-sensing-gis:5}
(a)~En unidades SI $L = \qty{9.0e6}{W.m^{-2}.sr^{-1}.m^{-1}}$. Con $c_1 =
2hc^2 = \qty{1.191e-16}{W.m^2.sr^{-1}}$, $c_2 = \qty{0.014388}{m.K}$ y
$\lambda^5 = (\num{1.09e-5})^5 = \num{1.539e-25}$, $c_1/(\lambda^5L) = 86.0$,
y $T_b = c_2/(\lambda\ln 87.0) = \qty{295.6}{K}$. (b)~La temperatura cinética
se obtiene de $0.98\,B_\lambda(T) = L$, es decir $B_\lambda(T) = L/0.98$: la misma
fórmula da $T = \qty{296.9}{K}$, \qty{1.3}{K} más.
\end{solution}

\begin{solution}{exo:b3:remote-sensing-gis:6}
(a)~$R = 693/\cos\ang{35} = \qty{846}{km}$. (b)~Antena real: $R\lambda/L =
\num{8.46e5}\times0.0555/12.3 = \qty{3.8}{km}$; apertura sintética: $L/2 =
\qty{6.2}{m}$, unas 600 veces más fina. (c)~La apertura sintética es la
longitud de la trayectoria a lo largo de la cual un punto permanece en el haz, $R\lambda/L =
\qty{3.8}{km}$.
\end{solution}

\begin{solution}{exo:b3:remote-sensing-gis:7}
$\partial z/\partial x = (418 - 405)/60 = 0.217$ (E menos O),
$\partial z/\partial y = (418 - 401)/60 = 0.283$ (N menos S);
$|\nabla z| = 0.357$, una pendiente de $\arctan 0.357 = \ang{19.6}$. El terreno
asciende hacia el nordeste, por lo que mira hacia el suroeste: $A =
\operatorname{atan2}(-0.217, -0.283) = \ang{217}$. Sombreado del relieve: $\cos\ang{45}
\cos\ang{19.6} + \sin\ang{45}\sin\ang{19.6}\cos(\ang{315} - \ang{217}) =
0.666 - 0.033 = 0.63$, algo más oscuro que el terreno llano (0.71): la pendiente
se aleja ligeramente de una fuente de luz situada al noroeste.
\end{solution}

\begin{solution}{exo:b3:remote-sensing-gis:8}
Totales de fila 52, 48, 38, 22; totales de columna 54, 46, 38, 22; diagonal 133 de
160: exactitud global \qty{83.1}{\percent}. Exactitudes del usuario: agua
$48/52 = 0.92$, bosque $37/48 = 0.77$, cultivos $29/38 = 0.76$, zona urbana
$19/22 = 0.86$. Del productor: agua $48/54 = 0.89$, bosque $37/46 = 0.80$,
cultivos $29/38 = 0.76$, zona urbana $19/22 = 0.86$. $p_e = (52\cdot54 + 48\cdot46 +
38\cdot38 + 22\cdot22)/160^2 = 0.271$, $\kappa = (0.831 - 0.271)/(1 -
0.271) = 0.77$. Los cultivos son los que con más frecuencia se omiten: 9 de los 38 píxeles reales de cultivo
fueron a parar a otras clases, sobre todo a bosque.
\end{solution}

\begin{solution}{exo:b3:remote-sensing-gis:9}
(a)~$L = 0.04\,B_\lambda(1443) + 0.96\,B_\lambda(290)$ a
\qty{10.9}{\micro\metre}, \qty{28.6}{W.m^{-2}.sr^{-1}.\micro\metre^{-1}},
cuya \omterm{def:b3:remote-sensing-gis:brightness}{temperatura de brillo} es \qty{396}{K}, frente a una temperatura media de
$0.04\times1443 + 0.96\times290 = \qty{336}{K}$. (b)~A
\qty{2.2}{\micro\metre}, $B_\lambda(1443)/B_\lambda(290) = \num{6.8e7}$:
la lava domina por completo el píxel, y su \omterm{def:b3:remote-sensing-gis:brightness}{temperatura de brillo} es
\qty{845}{K}. La banda del infrarrojo de onda corta detecta pequeños focos calientes mucho
mejor; la banda térmica, sensible a todo el píxel, es la que
mide el terreno que los rodea.
\end{solution}

\begin{solution}{exo:b3:remote-sensing-gis:10}
Con el satélite al este, $d_1 = u_E\sin\theta + u_z\cos\theta$; al
oeste, $d_2 = -u_E\sin\theta + u_z\cos\theta$. Sumando y restando,
$u_z = (d_1 + d_2)/(2\cos\theta) = 102/1.638 = \qty{62.3}{mm}$ y $u_E =
(d_1 - d_2)/(2\sin\theta) = 26/1.147 = \qty{22.7}{mm}$. Una pasada da una
ecuación para dos incógnitas: cualquier combinación de levantamiento y movimiento hacia el este con
la misma proyección sobre la \omterm{def:b3:remote-sensing-gis:insar}{línea de visión} da las mismas franjas.
\end{solution}

\begin{solution}{exo:b3:remote-sensing-gis:11}
(a)~Por el \cref{thm:b3:remote-sensing-gis:fringe}, $\Delta\phi =
4\pi\Delta r/\lambda < \pi$ exige que $\Delta r < \lambda/4 =
\qty{13.9}{mm}$. (b)~$u_z = k z_0(r^2 + z_0^2)^{-3/2}$ con $k = (1 -
\nu)\Delta V/\pi$; $\mathrm du_z/\mathrm dr = -3kz_0r(r^2 + z_0^2)^{-5/2}$,
cuya derivada se anula donde $(r^2 + z_0^2) = 5r^2$, en $r = z_0/2$. Allí
\[
\left|\frac{\mathrm du_z}{\mathrm dr}\right| = \frac{3k}{2z_0^3}\Big(\frac54\Big)^{-5/2}
= \frac{0.859\,k}{z_0^3} = \frac{0.859\times\num{2.39e6}}{\num{6.4e10}} = \num{3.2e-5},
\]
\qty{32}{mm} por kilómetro. (c)~$13.9/32 = \qty{0.43}{km}$: píxeles menores
que unos \qty{430}{m}, condición que cumplen fácilmente los de Sentinel-1.
\end{solution}

\begin{solution}{exo:b3:remote-sensing-gis:12}
(a)~Para $N, R \ge 0$, no ambos nulos, $|N - R| \le N + R$, luego $|\mathrm{NDVI}|
\le 1$; al sustituir $N$ y $R$ por $aN$ y $aR$ se multiplican el numerador y
el denominador por $a$. (b)~Sin corrección: $(0.502 - 0.051)/(0.502 + 0.051) =
0.816$; con corrección: $(0.512 - 0.081)/(0.512 + 0.081) = 0.727$. (c)~La radiancia
de trayectoria reduce el índice de la vegetación verde, en torno a una décima aquí, y
en una cantidad que cambia con la calima: un mapa de imágenes sin corregir mostraría
el tiempo atmosférico tanto como la vegetación, y las comparaciones entre fechas
carecerían de sentido.
\end{solution}

\begin{solution}{pb:b3:remote-sensing-gis:1}
\textbf{1.} Bosque $(0.38 - 0.04)/0.42 = 0.81$; colada de 1890 $0.23/0.37 =
0.62$; de 1960 $0.10/0.32 = 0.31$; de 2005 $0.02/0.26 = 0.08$.

\textbf{2.} 1890, luego 1960, luego 2005: la colada más antigua es la más verde,
en el orden de las edades; todavía no ha alcanzado el bosque.

\textbf{3.} La roca desnuda refleja algo más en el infrarrojo cercano que en
el rojo, como la arena parda de la figura del capítulo; los líquenes y las primeras
plantas añaden algo: 0.08 es el valor de un terreno casi desnudo.

\textbf{4.} Una ladera orientada en sentido contrario a la luz atenúa ambas bandas en la misma
proporción; el cociente lo cancela, mientras que la reflectancia en el infrarrojo cercano por sí sola
confundiría un bosque sombreado con vegetación dispersa.

\textbf{5.} Mediante la \cref{prop:b3:remote-sensing-gis:planck-inversion} con
$\lambda = \qty{10.9}{\micro\metre}$: la meseta \qty{291.1}{K}, el
píxel de la fumarola \qty{302.3}{K}.

\textbf{6.} $0.97\,B_\lambda(T) = \num{8.40e6}$ da $T = \qty{293.0}{K}$.

\textbf{7.} $B_\lambda(\qty{370}{K}) =
\qty{22.48}{W.m^{-2}.sr^{-1}.\micro\metre^{-1}}$; $f = (9.95 - 8.40)/(22.48
- 8.40) = 0.110$, unos \qty{1100}{m^2} del píxel de \qty{e4}{m^2}.

\textbf{8.} La radiancia es una función convexa de la temperatura
(véase la \cref{prop:b3:remote-sensing-gis:mixed}): la parte caliente aporta más radiancia
que su parte proporcional de la temperatura, de modo que el píxel parece más caliente que su media,
$0.110\times370 + 0.890\times291.1 = \qty{299.8}{K}$ frente a
\qty{302.3}{K}.

\textbf{9.} La fase de cada eco depende de la trayectoria exacta; un cambio grande
de punto de vista cambia la suma de los ecos de los dispersores y el
\omterm{def:b3:remote-sensing-gis:insar}{interferograma} pierde su coherencia, y además añade una fase debida al relieve.
$48/12 = 4$ ciclos de un solo satélite.

\textbf{10.} $5.4\times27.73 = \qty{149.7}{mm}$ hacia el satélite.

\textbf{11.} $u_z = 149.7/\cos\ang{35} = \qty{182.8}{mm}$.

\textbf{12.} El terreno se desplaza hacia fuera además de hacia arriba; en el lado este el
movimiento hacia fuera va hacia el satélite y se suma al levantamiento; en el lado oeste
lo resta (véase la \cref{prop:b3:remote-sensing-gis:mogi-los}). El cambio
cae más deprisa hacia el oeste, donde las franjas están apretadas; el patrón
se estira hacia el este.

\textbf{13.} En esa línea el movimiento horizontal es norte--sur,
perpendicular a la dirección de observación, de modo que el radar ve solo $u_z\cos\theta$, con la
forma del levantamiento. El levantamiento se reduce a la mitad en $r = 0.766\,z_0$, así que $z_0 =
3.4/0.766 = \qty{4.44}{km}$.

\textbf{14.} La mitad de 5.4: 2.7 franjas.

\textbf{15.} Una deformación es igual en los dos \omterm{def:b3:remote-sensing-gis:insar}{interferogramas} tras
normalizar por los intervalos de tiempo, o sigue creciendo de un par a otro; un
patrón atmosférico cambia de una fecha a otra y a menudo sigue el
relieve. Un segundo par que muestre la misma diana confirma la
deformación.

\textbf{16.} $\Delta V = \pi u_z(0)z_0^2/(1 - \nu) = \pi\times0.1828\times
4436^2/0.75 = \qty{1.51e7}{m^3}$.

\textbf{17.} $\num{1.51e7}/48 = \qty{3.1e5}{m^3}$ al día, \qty{1.1e8}{m^3}
al año si se mantuviera el ritmo.

\textbf{18.} $u_r \propto r/(r^2 + z_0^2)^{3/2}$ es máximo donde su
derivada se anula, $r^2 + z_0^2 = 3r^2$, $r = z_0/\sqrt2 = \qty{3.14}{km}$;
allí $u_r = \qty{70}{mm}$, \qty{38}{\percent} del levantamiento central.

\textbf{19.} Radio \qty{1.48}{km}, volumen $\frac43\pi(1.48)^3 =
\qty{13.6}{km^3}$; el cambio es $\num{1.51e7}/\num{1.36e10} =
\qty{0.11}{\percent}$ de este.

\textbf{20.} No: una décima de franja es el \qty{2}{\percent} del recuento, y
la profundidad, leída a partir de una semianchura de \qty{3.4}{km}, quizá hasta
\qty{0.1}{km}, entra al cuadrado, \qty{6}{\percent}; la corrección del
\qty{3}{\percent} queda dentro de esa incertidumbre.

\textbf{21.} El calentamiento del campo de fumarolas (parte~II) a lo largo del tiempo, un
aumento de sismos pequeños bajo el abombamiento, un cambio en el gas
liberado, y el levantamiento medido mediante posicionamiento por satélite o nivelación en
unos pocos puntos del flanco.

\textbf{22.} El reservorio bajo el monte Halvren creció en
\textbf{15~millones~\unit[text-series-to-math=true]{m^3}} ($\num{1.51e7}\,\unit{m^3}$).
\end{solution}
